%!TEX root = ../main.tex

\begin{abstract}[it]

Il question answering conversazionale su collezioni di documenti pone sfide
specifiche rispetto al question answering su dominio aperto: le risposte devono
essere radicate in un corpus definito, le domande sono spesso decontestualizzate
nel corso di più turni di dialogo, e alcune domande non trovano risposta nei
documenti disponibili.

Questa tesi presenta una valutazione empirica sistematica di approcci a questo
problema, utilizzando \gls{mdd} come benchmark principale --- un dataset di
dialoghi orientati agli obiettivi, fondati su documenti di servizi governativi.
Partendo da un sistema basato esclusivamente su \gls{llm} (\fone{} = 15,9\% sul
validation set completo), viene sviluppata progressivamente una pipeline completa
di \gls{rag} con recupero ibrido sparso--denso e riscrittura conversazionale delle
query, raggiungendo un \fone{} di 26,8\%.

Un esperimento con contesto oracle (sezioni di riferimento fornite direttamente al
modello) stabilisce un tetto di generazione del 36,1\% di \fone{}; circa due terzi
del divario rispetto al miglior sistema \gls{rag} sono dovuti a errori di recupero.
Un reranker cross-encoder migliora la precisione del recupero (\ratk{5}: da 61,2\%
a 67,4\%) senza produrre un guadagno statisticamente significativo in \fone{}:
l'analisi per esempio mostra che il generatore sfrutta pienamente le nuove sezioni
pertinenti recuperate, ma il guadagno è in gran parte compensato dalle sezioni
corrette che il reranker esclude dalle prime cinque.

La tesi si conclude con una descrizione delle attività future pianificate:
costruzione e interrogazione di grafi della conoscenza, architetture Graph-\gls{rag}
e approcci basati su agenti.

\end{abstract}
